% À minha mãe, \textit{Graciela}, cuja luta diária foi e sempre será minha maior fonte de inspiração. Obrigado por me ensinar a persistir. Às minhas irmãs, que carrego nos pensamentos mais emotivos e profundos que consigo conceber; vocês são parte essencial de quem sou.

% Um agradecimento e adeus saudoso aos meus companheiros de casa, no Peru, que partiram para outra vida enquanto eu estava longe. À \textit{Wanda}, minha cadela levada por um acidente; à \textit{Chuleta} e \textit{Luna}, minhas gatas que partiram sob circunstâncias que não mereciam; e à \textit{Preciosa}, a cadelinha de quem guardo inúmeras memórias e que partiu com o peso dos anos.

% Aos amigos peruanos e brasileiros que fiz nesta jornada. Agradeço a vocês por terem cruzado meu caminho, muitas vezes mesmo quando eu tentava evitar ao máximo interações sociais. Obrigado por existirem e por me provarem, contra a minha própria resistência, que sempre haverá pessoas interessantes a serem conhecidas.

% Aos camaradas do laboratório MASTER, por todos os dias e noites, dias de semana e finais de semana que passamos aqui. Mesmo como "bichos do mato", vocês me ajudaram de formas diferentes: com as gírias, com as palavras e suportando minhas dúvidas e minha audição baixa, que dificultava entender o português que vocês falam e eu aprendo. Em especial a Danilo e Jean, que foram vítimas das minhas dúvidas sobre a língua — muitas vezes sem sentido para vocês, mas que para mim faziam muita diferença.

% Ao meu orientador, Prof. Dr. Ricardo Camargo, não apenas pela orientação acadêmica, mas por aceitar trabalhar comigo. Obrigado pelos conselhos que me ajudaram a navegar e me adaptar a esta sociedade brasileira — um lugar que, curiosamente, se revelou não tão diferente da minha origem, mas ao mesmo tempo, profundamente distinto em suas nuances.

Agradeço aos membros da banca examinadora pela aceitação do convite para avaliar meu trabalho, contribuindo para a melhoria do mesmo e para meu desenvolvimento profissional e acadêmico.

À Petrobras e à FDTE (Fundação para o Desenvolvimento Tecnológico da Engenharia), pelo suporte financeiro essencial através da bolsa de estudos, que permitiu a realização deste trabalho. Sua existência incentivam o avanço da ciência, tecnologia e inovação em Brasil, mantendo a pesquisa científica viva e ativa.

Ao Instituto de Astronomia, Geofísica e Ciências Atmosféricas (IAG) e à Universidade de São Paulo (USP), pela infraestrutura e excelência acadêmica.

\vfill

\begin{flushleft}
\rule{6cm}{0.5pt}\\
{\footnotesize{Esta tese/dissertação foi escrita em \LaTeX{} com a classe IAGTESE, para teses e dissertações do IAG.}}
\end{flushleft}